\frfilename{mt21.tex} 
\versiondate{17.1.15} 
\copyrightdate{1994} 
 
\def\chaptername{Taksonomi ruang ukur} 
 
\gdef\topparagraph{} 
   \gdef\bottomparagraph{Bab\ 21 {\it pendahuluan.}} 
   \gdef\newparagraph{Bab\ 21 {\it pendahuluan.}} 
   \gdef\sectionname{Pendahuluan} 
   \gdef\headlinesectionname{Pendahuluan} 
   \centerline{\bf *Bab 21} 
   \medskip 
   \centerline{\bf \chaptername} 
   \medskip 
 
Saya memulai jilid ini dengan sebuah `bab berbintang'. Maksudnya ialah
bahwa saya sebenarnya tidak menganjurkan bab ini bagi pemula. Bab ini
membahas beragam persoalan teknis yang sangat penting bagi perkembangan
selanjutnya dari pokok bahasan ini, tetapi mungkin terasa abstrak
sekaligus samar bagi siapa pun yang belum menjumpai persoalan-persoalan
yang hendak diselesaikan oleh teknik-teknik tersebut. Di pihak lain,
jika (sebagaimana lazimnya) pembahasan ini ditiadakan dan gagasannya
baru diperkenalkan ketika benar-benar diperlukan, mahasiswa mungkin
pada akhirnya hanya memiliki gambaran yang sangat kabur mengenai
teorema mana yang dapat diharapkan berlaku bagi jenis ruang ukur yang
mana, serta tanpa kosakata untuk mengungkapkan gagasan-gagasan tersebut.
Karena itu, saya menggunakan beberapa halaman untuk memperkenalkan
terminologi dan konsep yang dapat dipakai untuk membedakan ruang ukur
yang `baik' dari ruang ukur lainnya, disertai beberapa hubungan dasarnya.
Paragraf-paragraf yang langsung relevan dengan teori yang dipaparkan
dalam Jilid 1 hanyalah paragraf tentang ruang ukur `lengkap',
`$\sigma$-hingga', dan `semihingga' (211A, 211D, 211F, 211Lc, \S212,
213A-213B, 215B), serta tentang ukuran Lebesgue (211M). Untuk selebihnya,
menurut saya seorang pendatang baru dalam bidang ini dapat melewati bab
ini untuk sementara waktu, lalu kembali ke bagian-bagian tertentu
ketika teks dalam bab-bab berikutnya merujuk kepadanya.
Di pihak lain, bab ini juga dapat digunakan sebagai pengantar pada corak
teori ukuran yang `paling murni', yakni kajian ruang ukur sebagai
objek tersendiri, dengan pembahasan sistematis atas beberapa
konstruksi elementer.
 
\discrpage 




